\documentclass[../note.tex]{subfiles}
\usepackage{bm}

\begin{document}


\chapter{Bayesian Linear Regression}

\section{Basic Settings}

Consider a training dataset 
\[
    \mathcal{D} = \left\{(x_i,y_i)\right\}_{i=1}^{N},
\]
consisting of $N$ observations, where $\bm{x}_i \in \mathbb{R}^{p}$ denotes the $p$-dimensional feature vector associated with the $i$-th observation and $y_i \in \mathbb{R}$ denotes the corresponding scalar response. 

Define design matrix $\mathbf{X}$, target vector $\bm{y}$, and coefficient vector $\bm{w}$ as 
\[
    \mathbf{X} = \begin{bmatrix}
        \bm{x}_1^{\top} \\
        \vdots \\
        \bm{x}_N^{\top}
    \end{bmatrix} \in \mathbb{R}^{N \times p}, \qquad
    \bm{y} = \begin{bmatrix}
        {y}_1 \\
        \vdots \\
        {y}_N
    \end{bmatrix} \in \mathbb{R}^{N}, \qquad
    \bm{w} = \begin{bmatrix}
        {w}_1 \\
        \vdots \\
        {w}_p
    \end{bmatrix} \in \mathbb{R}^{p}.
\]

Under the linear regression model, each observation satisfies
\[
    y_i = \bm{w}^{\top} \bm{x}_i + \epsilon_i,
\]
where the noise terms are assumed to be independent and identically distributed as
\[
    \epsilon_i \overset{\mathrm{i.i.d.}}{\sim}\mathcal N(0,\sigma^2), \qquad i=1,\dots,N.
\]

Define the noise vector
\[
    \bm{\epsilon} = \begin{bmatrix}
        \epsilon_1 \\
        \vdots \\
        \epsilon_N
    \end{bmatrix}.
\]
It follows that
\[
    \bm\epsilon \sim \mathcal N(\bm 0,\sigma^2\mathbf I_N),
\]
where $\mathbf{I}_N$ denotes the $N\times N$ identity matrix.

For the entire dataset, the linear regression model can therefore be written as
\[
    \bm{y} = \mathbf{X}\bm{w} + \bm{\epsilon}.
\]
Consequently,
\[
    p(\bm{y} | \mathbf{X}, \bm{w}) \sim \mathcal{N}(\mathbf{X}\bm{w}, \sigma \mathbf{I}_N).
\]

Bayesian method includes two basic steps: inference and prediction, where inference is to estimate posterior $\bm{w}$, while prediction is to predict ${y}^{*}$ given $\bm{x}^{*}$.

\section{Inference}


\subsection{Model Establishment}

In Bayesian linear regression, the coefficient vector $\bm{w}$ is treated as a random variable and assigned a prior distribution. We assume that the prior distribution of $\bm{w}$ does not depend on $\mathbf{X}$, i.e., $w\perp X$, and thus $p(\bm w\mid\mathbf X)=p(\bm w)$.

Given the training dataset, Bayesian formula is decomposed as: 
\[
    p(\bm{w} | \mathcal{D}) = p(\bm{w} | \mathbf{X}, \bm{y}) = \frac{p(\bm{w}, \mathbf{X}, \bm{y})}{\textcolor{red}{p(\mathbf{X}, \bm{y})}} = \frac{\textcolor{red}{p(\bm{w}, \mathbf{X}, \bm{y})}}{p(\bm{y}|\mathbf{X}) \textcolor{red}{p(\mathbf{X})} } = \frac{\textcolor{red}{p(\bm{w}, \bm{y} | \mathbf{X})}}{p(\bm{y}|\mathbf{X})} = \frac{p(\bm{y} | \mathbf{X}, \bm{w}) \textcolor{red}{p(\bm{w} | \mathbf{X})}}{p(\bm{y}|\mathbf{X})} = \frac{p(\bm{y} | \mathbf{X}, \bm{w}) p(\bm{w})}{p(\bm{y}|\mathbf{X})},
\]
where the fifth equality follows from
\[
    p(\bm{y} | \mathbf{X}, \bm{w}) = \frac{p(\bm{w}, \mathbf{X}, \bm{y})}{p(\bm{w}, \mathbf{X})} = \frac{p(\bm{w}, \mathbf{X}, \bm{y})}{ p(\bm{w} | \mathbf{X}) p(\mathbf{X})} = \frac{p(\bm{w}, \bm{y} | \mathbf{X})}{p(\bm{w} | \mathbf{X})},
\]
and the last equality follows from independency of prior coefficient and design matrix. 

The denominator 
\[
    p(\bm{y}\mid\mathbf{X}) = \int p(\bm{y}\mid\mathbf{X},\bm{w}) p(\bm{w})\mathrm{d}\bm{w}
\]  
is the marginal likelihood. Since it does not depend on $\bm{w}$, the posterior distribution can be expressed up to a proportionality constant as 
\[
    p(\bm{w} | \mathbf{X}, \bm{y}) \propto p(\bm{y} | \mathbf{X}, \bm{w}) p(\bm{w}),
\]
where $p(\bm{y} | \mathbf{X}, \bm{w})$ is a likelihood function and $p(\bm{w})$ is a prior function.

Suppose that a zero-mean Gaussian prior is imposed on the coefficient vector:
\[
    p(\bm{w}) = \mathcal{N}\left(\bm{0}, \mathbf{\Sigma}_{p}\right),
\]
where $\mathbf{\Sigma}_p\in\mathbb{R}^{p\times p}$ denotes the prior covariance matrix.

Since both the likelihood and the prior are Gaussian, their conjugacy implies that the posterior distribution of $\bm{w}$ is also Gaussian:
\[
    p(\bm{w} | \mathbf{X}, \bm{y}) = \mathcal{N}(\mu_w, \Sigma_w)
\]
Thus,
\[
    p(\bm{w} | \mathbf{X}, \bm{y}) \propto \mathcal{N}(\mathbf{X}\bm{w}, \sigma \mathbf{I}_N) \mathcal{N}\left(\bm{0}, \mathbf{\Sigma}_{p}\right).
\]

The remaining task is therefore to derive the posterior mean $\bm{\mu}_w$ and posterior covariance $\mathbf{\Sigma}_w$.

\subsection{Model Solution}



\section{Prediction}

\end{document}